\chapter{Hash Tables}

Chapter~10 established the fundamental trade-off of elementary data structures:
contiguous arrays offer instant $\Theta(1)$ indexing but expensive $\Theta(n)$
insertions, while linked lists offer cheap $\Theta(1)$ pointer splices but force
$\Theta(n)$ linear scans to find an element by key. Neither structure can search
for an arbitrary key in constant time.

Imagine you run a small library, and someone asks for book number 4173. You could
walk the shelves one by one, inspecting every spine -- that is searching a linked
list. Or you could build the library so that book 4173 lives at position 4173 on a
very long shelf, letting you walk straight to it using array index arithmetic. The
catch is that a shelf with a slot for every possible catalog number wastes immense
space if most numbers never appear. Hash tables deliver the best of both worlds:
they compute memory locations directly from keys to achieve $\Theta(1)$ average
search, insertion, and deletion, while compressing a vast universe of possible keys
into an array proportional to the records actually stored.

By the end of this chapter you will understand why a well-built hash table answers ``is $x$ in here?'', ``insert $x$'', and ``delete $x$'' in \emph{constant expected time}, what ``expected'' is doing in that sentence, and how to build hash functions that earn that guarantee instead of just hoping for it.

\begin{focusbox}
A hash table is a dictionary: it stores a set of \emph{(key, value)} pairs and supports \textsc{Insert}, \textsc{Search}, and \textsc{Delete}. Its trick is a \emph{hash function} $h$ that maps each key into a slot of a small array. The core architectural challenge is scattering keys so uniformly across the array that every slot holds only a tiny handful of elements, keeping lookup times virtually instantaneous while using memory proportional only to the keys actually stored.
\end{focusbox}

\section{The dictionary problem}

First, let us specify the exact operations we need. We have a universe $U$ of possible keys -- maybe all 32-bit integers, maybe all strings up to length 100, maybe all possible employee IDs. At any moment we are storing some \emph{subset} of those keys, a dynamic set we will call $K$, and $K$ is usually \emph{tiny} compared to $U$. We want three operations:

\begin{itemize}
  \item \textsc{Search}$(T, k)$: return the element with key $k$, or report that it is absent.
  \item \textsc{Insert}$(T, x)$: add element $x$ (whose key is $x.key$) to the set.
  \item \textsc{Delete}$(T, x)$: remove element $x$ from the set.
\end{itemize}

We do \emph{not} need the keys sorted, we do \emph{not} need to ask ``what is the smallest key larger than $k$'', we do \emph{not} need predecessors or successors. We just need membership and retrieval. That modest wish list is exactly why hash tables can be so fast: by giving up order, they buy speed.

\begin{aestheticbox}
This is one of the most-used data structures on Earth. Python's \lstinline|dict| and \lstinline|set|, Java's \lstinline|HashMap|, C++'s \lstinline|unordered_map|, JavaScript objects, Ruby hashes, the symbol tables inside every compiler, the routing tables inside databases, deduplication, caching, and the ``have I seen this before?'' check in a hundred algorithms -- all hash tables underneath. When people say a program ``uses a hash map to make it fast,'' they mean: it turned an $O(n)$ scan into an $O(1)$ lookup.
\end{aestheticbox}

\section{Direct-address tables: the simple, wasteful idea}

Let us start with the version that has no cleverness at all, just to see the shape of the answer before we economize.

Suppose the universe of keys is small: $U = \{0, 1, 2, \dots, m-1\}$ for some manageable $m$, and no two elements share a key. Then we can do the library-shelf thing literally. Allocate an array $T[0\,..\,m-1]$, called a \emph{direct-address table}. Slot $T[k]$ holds the element whose key is $k$ (or a pointer to it), and holds \textsc{nil} if no such element is present.

\begin{algorithm}[H]
\caption{Direct-address operations}
\KwIn{A direct-address table $T[0\,..\,m-1]$, a key $k$, an element $x$}
\textsc{Direct-Address-Search}$(T,k)$: \Return $T[k]$\;
\textsc{Direct-Address-Insert}$(T,x)$: $T[x.key] \gets x$\;
\textsc{Direct-Address-Delete}$(T,x)$: $T[x.key] \gets \textsc{nil}$\;
\end{algorithm}

Each operation is three lines, no loops, no comparisons. Every one runs in $O(1)$ time in the \emph{worst} case -- not expected, not amortized, just flatly constant. This is the dream we are chasing.

\begin{example}
Say employee IDs run from $0$ to $99$, and right now we employ people with IDs $4$, $17$, and $42$. The direct-address table is an array of $100$ slots. $T[4]$, $T[17]$, $T[42]$ point at their records; the other $97$ slots hold \textsc{nil}. ``Is employee $17$ here?'' is just ``look at $T[17]$'' -- one memory access.
\end{example}

The catch is memory: when the universe of possible keys $U$ is large, allocating one slot per possible key is impractical or impossible.

\begin{trapbox}
\textbf{Direct addressing wastes space when $U$ is large.} If keys are 32-bit integers, $U$ has about $4.3$ billion values, so $T$ needs $4.3$ billion slots even to store ten keys. That is gigabytes of mostly-\textsc{nil} array. And if keys are strings or 64-bit numbers, a direct-address table is not merely wasteful -- it is physically impossible to allocate. The whole point of hashing is to escape this.
\end{trapbox}

\begin{intuitionbox}
Direct addressing pays $|U|$ memory to get $O(1)$ time. We are about to pay only $\Theta(|K|)$ memory -- proportional to the number of keys we actually store -- and \emph{keep} $O(1)$ time, on average. The trade is that ``on average'' replaces ``always.'' For almost every real program that is a fantastic bargain.
\end{intuitionbox}

\section{From direct addressing to hashing}

Here is the leap. We have a huge universe $U$ but only ever store $|K| = n$ keys, and we are willing to allocate an array $T[0\,..\,m-1]$ of just $m$ slots, where $m$ is on the order of $n$, not $|U|$. The problem: a key $k$ might be a nine-digit number, but our array only has, say, $m = 701$ slots. Where does key $k$ go?

We \emph{compute} a slot from the key, using a \emph{hash function}
\[
  h : U \to \{0, 1, \dots, m-1\}.
\]
The key $k$ ``hashes to'' slot $h(k)$, and we say $h(k)$ is the \emph{hash value} of $k$. Instead of storing element $x$ at $T[x.key]$ (impossible -- the array is too small), we store it at $T[h(x.key)]$. The array shrinks from $|U|$ slots to $m$ slots.

But you already feel the snag. We are cramming $|U|$ possible keys into $m \ll |U|$ slots. By the pigeonhole principle, \emph{many keys must map to the same slot}. When $h(k_1) = h(k_2)$ for $k_1 \neq k_2$, we call it a \emph{collision}, and we now have two elements wanting the same parking spot.

Why do collisions happen so much faster than intuition suggests? If you have $m = 365$ slots and begin inserting keys, naive intuition might suggest that collisions remain rare until the table is mostly full. Yet, by the classic birthday paradox, the probability of at least one collision crosses $50\%$ after inserting just $n \approx \sqrt{2m \ln 2} \approx 23$ keys. Because the number of possible colliding pairs grows quadratically as $\binom{n}{2} = n(n-1)/2$, collisions begin scattering across slots long before the table approaches capacity. Managing collisions is therefore the central structural pillar of hash table design, not an afterthought.

\begin{figure}[h]
\centering
\begin{tikzpicture}[font=\small]
  % universe blob
  \draw[rounded corners=14pt, fill=sourceblue!8, draw=sourceblue!70]
    (-0.2,-2.6) rectangle (3.2,0.8);
  \node[sourceblue!60!black] at (1.5,0.45) {Universe $U$ (huge)};
  \node[draw,circle,fill=white,inner sep=1.2pt] (k1) at (0.7,-0.4) {$k_1$};
  \node[draw,circle,fill=white,inner sep=1.2pt] (k2) at (2.3,-0.8) {$k_2$};
  \node[draw,circle,fill=white,inner sep=1.2pt] (k3) at (0.9,-1.7) {$k_3$};
  \node[draw,circle,fill=white,inner sep=1.2pt] (k4) at (2.4,-1.9) {$k_4$};

  % table slots
  \foreach \i in {0,...,5} {
    \draw (6,-0.5*\i) rectangle (7.4,-0.5*\i+0.5);
    \node[gray!70, font=\scriptsize] at (6.7,-0.5*\i+0.25) {[\i]};
  }
  \node[font=\small\bfseries] at (6.7,0.75) {Table $T[0\,..\,5]$};

  % arrows
  \draw[-{Latex},sourceblue!70] (k1) to[out=0,in=180] (6,0.25);
  \draw[-{Latex},sourceblue!70] (k3) to[out=0,in=180] (6,0.25);
  \draw[-{Latex},targetgreen!80] (k2) to[out=0,in=180] (6,-1.25);
  \draw[-{Latex},targetgreen!80] (k4) to[out=0,in=180] (6,-2.25);
  \node[trapred, font=\scriptsize, anchor=west] at (7.7,0.25) {$h(k_1)=h(k_3)$: collision!};
\end{tikzpicture}
\caption{A hash function squashes a huge universe into a small table. Keys $k_1$ and $k_3$ collide in slot $0$.}
\label{ch11:hashing-picture}
\end{figure}

\section{Collision resolution by chaining}

The simplest, most robust answer to collisions is: if several keys land in the same slot, just \emph{keep a list there}. Slot $T[j]$ holds not a single element but a pointer to a linked list of all elements whose keys hash to $j$. This is \emph{chaining}.

\begin{intuitionbox}
Picture a coat-check with $m$ hooks. The hash function tells you which hook your coat goes on. If two coats get the same hook, you just hang the second one in front of the first -- the hook holds a little chain of coats. To find your coat, you go to its hook and look through the (hopefully short) chain hanging there.
\end{intuitionbox}

The three operations become:

\begin{algorithm}[H]
\caption{Chained hash table operations}
\KwIn{Hash table $T$ with $m$ slots, hash function $h$, element $x$, key $k$}
\textsc{Chained-Hash-Insert}$(T,x)$: insert $x$ at the head of list $T[h(x.key)]$\;
\textsc{Chained-Hash-Search}$(T,k)$: search list $T[h(k)]$ for an element with key $k$\;
\textsc{Chained-Hash-Delete}$(T,x)$: delete $x$ from list $T[h(x.key)]$\;
\end{algorithm}

\textsc{Insert} is $O(1)$ worst case: hash the key, splice the new node onto the front of the list -- no scanning, because we assume $x$ is not already present (or we do not care if it duplicates). \textsc{Search} must walk the list at slot $h(k)$, so its cost is proportional to that list's length. \textsc{Delete}, if we are handed a pointer to $x$ and use doubly linked lists, is also $O(1)$: we splice $x$ out without searching for it.

\begin{figure}[h]
\centering
\begin{tikzpicture}[font=\small]
  % Table header and slots
  \node[font=\footnotesize\scshape, text=gray!70!black] at (0.5, 0.9) {Slots};
  \node[font=\footnotesize\scshape, text=gray!70!black] at (3.5, 0.9) {Chained Buckets};

  \foreach \i in {0,...,4} {
    \draw[thick, draw=gray!60, fill=gray!8, rounded corners=1pt] (0, -0.8*\i) rectangle (1.1, -0.8*\i+0.6);
    \node[gray!70!black, font=\scriptsize] at (-0.4, -0.8*\i+0.3) {[\i]};
    \fill[gray!50] (0.55, -0.8*\i+0.3) circle (2pt);
  }

  % Slot 0: empty
  \draw[-{Latex[length=2.2mm]}, thick, gray!60] (0.55, 0.3) -- (1.6, 0.3);
  \node[gray!60, font=\scriptsize\sffamily] at (1.9, 0.3) {nil};

  % Slot 1: two elements (19 -> 5)
  \draw[-{Latex[length=2.2mm]}, thick, targetgreen!80] (0.55, -0.5) -- (1.5, -0.5);
  \node[draw=targetgreen!80, fill=targetgreen!12, rounded corners=2pt, inner sep=4pt] (a1) at (2.4, -0.5) {$k{=}19$};
  \draw[-{Latex[length=2.2mm]}, thick, targetgreen!80] (a1.east) -- ++(0.6,0) node[draw=targetgreen!80, fill=targetgreen!12, rounded corners=2pt, inner sep=4pt, anchor=west] (a2) {$k{=}5$};
  \draw[-{Latex[length=2.2mm]}, thick, targetgreen!80] (a2.east) -- ++(0.5,0) node[gray!60, font=\scriptsize\sffamily, anchor=west] {nil};

  % Slot 2: empty
  \draw[-{Latex[length=2.2mm]}, thick, gray!60] (0.55, -1.3) -- (1.6, -1.3);
  \node[gray!60, font=\scriptsize\sffamily] at (1.9, -1.3) {nil};

  % Slot 3: one element (17)
  \draw[-{Latex[length=2.2mm]}, thick, quotientorange!80] (0.55, -2.1) -- (1.5, -2.1);
  \node[draw=quotientorange!80, fill=quotientorange!15, rounded corners=2pt, inner sep=4pt] (b1) at (2.4, -2.1) {$k{=}17$};
  \draw[-{Latex[length=2.2mm]}, thick, quotientorange!80] (b1.east) -- ++(0.5,0) node[gray!60, font=\scriptsize\sffamily, anchor=west] {nil};

  % Slot 4: three elements (33 -> 12 -> 68) with search trace for k=12
  \draw[-{Latex[length=2.2mm]}, thick, sourceblue!80] (0.55, -2.9) -- (1.5, -2.9);
  \node[draw=sourceblue!80, fill=sourceblue!12, rounded corners=2pt, inner sep=4pt] (c1) at (2.4, -2.9) {$k{=}33$};
  \draw[-{Latex[length=2.2mm]}, thick, sourceblue!80] (c1.east) -- ++(0.6,0) node[draw=sourceblue, fill=sourceblue!25, thick, rounded corners=2pt, inner sep=4pt, anchor=west] (c2) {$k{=}12$};
  \draw[-{Latex[length=2.2mm]}, thick, sourceblue!80] (c2.east) -- ++(0.6,0) node[draw=sourceblue!80, fill=sourceblue!12, rounded corners=2pt, inner sep=4pt, anchor=west] (c3) {$k{=}68$};
  \draw[-{Latex[length=2.2mm]}, thick, sourceblue!80] (c3.east) -- ++(0.5,0) node[gray!60, font=\scriptsize\sffamily, anchor=west] {nil};

  % Annotation for search path
  \node[font=\scriptsize, text=trapred] (s1) at (2.2, -3.7) {step 1: $33 \neq 12$};
  \draw[-{Latex[length=2mm]}, thick, dashed, trapred] (s1.north) -- (c1.south);

  \node[font=\scriptsize, text=targetgreen!80!black] (s2) at (4.9, -3.7) {step 2: $12 = 12$ (found!)};
  \draw[-{Latex[length=2mm]}, thick, dashed, targetgreen!80!black] (s2.north) -- (c2.south);
\end{tikzpicture}
\caption{Collision resolution by chaining with $m=5$ slots. Keys $19$ and $5$ map to slot $1$; keys $33$, $12$, and $68$ map to slot $4$. Searching for $k=12$ computes $h(12)=4$, jumps to list head $T[4]$, and walks $2$ nodes before terminating successfully.}
\label{ch11:chaining-picture}
\end{figure}

\subsection{The load factor}

How long are these lists, on average? If the table has $m$ slots and currently holds $n$ elements, define the \emph{load factor}
\[
  \alpha = \frac{n}{m}.
\]
It is the average number of elements per slot. If $n = 1000$ and $m = 500$, then $\alpha = 2$: on average two elements per list. A load factor of $\alpha = 0.75$ means the table is about three-quarters full in the ``one element per slot'' sense. The load factor is the single most important dial on a hash table, and a great deal of our analysis is really analysis of $\alpha$.

\begin{focusbox}
$\alpha = n/m$ is the average list length under chaining. We will see that search costs $\Theta(1 + \alpha)$. So if we keep $\alpha$ bounded by a constant -- say, never let it exceed $1$ or $2$ -- searches are $\Theta(1)$. Keeping $\alpha$ small is why hash tables \emph{resize} (rehash into a bigger array) as they fill.
\end{focusbox}

\section{Analyzing chaining: simple uniform hashing}

We want to prove that search is fast on average. But ``on average'' over what? There is no randomness in a fixed set of keys and a fixed hash function -- if an adversary hands you keys that all collide, every search is slow, period. To get a meaningful average we make an idealizing assumption about \emph{how keys distribute}.

\begin{formalbox}
\textbf{Simple uniform hashing.} Any given key is equally likely to hash into any of the $m$ slots, independently of where the other keys have hashed. Formally, for the keys we store, $\Pr[h(k) = j] = 1/m$ for every slot $j$, and the slots chosen by distinct keys are independent.
\end{formalbox}

This is an assumption about the \emph{behavior} of $h$ on our particular keys, not something we can prove about an arbitrary function. Later, with \emph{universal hashing}, we will replace this leap of faith with a real probabilistic guarantee. For now, accept it and see where it leads.

Let $n_j$ denote the length of the list at slot $j$, so $n = n_0 + n_1 + \cdots + n_{m-1}$. Under simple uniform hashing, each of the $n$ keys lands in slot $j$ with probability $1/m$, so the \emph{expected} length of any list is
\[
  \mathbb{E}[n_j] = \frac{n}{m} = \alpha.
\]

\subsection{Unsuccessful search}

An \emph{unsuccessful} search -- looking for a key $k$ that is \emph{not} in the table -- walks the entire list at slot $h(k)$ to the end (it has to, to be sure $k$ is absent). Under simple uniform hashing, $k$ is equally likely to hash to any slot, and the expected length of that slot's list is $\alpha$. Add the constant cost of computing $h(k)$ and indexing into the array, and we get expected time
\[
  \Theta(1 + \alpha).
\]

\subsection{Successful search}

A \emph{successful} search is a touch subtler, because the key we are looking for is somewhere in its list, and on average we walk only \emph{partway} down. The careful count -- which we do gently below -- also gives $\Theta(1 + \alpha)$.

\begin{theorem}[Expected cost of search under chaining]
\label{ch11:chaining-theorem}
In a hash table with chaining, under the assumption of simple uniform hashing, both a successful and an unsuccessful search take expected time $\Theta(1 + \alpha)$, where $\alpha = n/m$ is the load factor.
\end{theorem}

\begin{proof}[Sketch for successful search]
Suppose we search for each of the $n$ stored keys equally often. Number the keys $x_1, x_2, \dots, x_n$ in the order they were \emph{inserted} (each at the head of its list). When we successfully search for $x_i$, the number of elements we examine is $1$ (for $x_i$ itself) plus the number of elements that appear \emph{before} $x_i$ in its list -- which are exactly the keys inserted \emph{after} $x_i$ that happened to land in the same slot.

For $j > i$, let the indicator $X_{ij} = 1$ if $x_i$ and $x_j$ hash to the same slot. Under simple uniform hashing, $\Pr[X_{ij}=1] = 1/m$, so $\mathbb{E}[X_{ij}] = 1/m$. The expected number of elements examined when searching for $x_i$ is $1 + \sum_{j>i} \mathbb{E}[X_{ij}]$. Averaging over all $i$,
\[
  \frac{1}{n}\sum_{i=1}^{n}\Big(1 + \sum_{j=i+1}^{n}\frac{1}{m}\Big)
  = 1 + \frac{1}{nm}\sum_{i=1}^{n}(n-i)
  = 1 + \frac{1}{nm}\cdot\frac{n(n-1)}{2}
  = 1 + \frac{n-1}{2m}.
\]
Since $\frac{n-1}{2m} < \frac{n}{2m} = \frac{\alpha}{2}$, the expected number of elements examined is at most $1 + \alpha/2$, giving total expected time $\Theta(1+\alpha)$ once we add the constant cost of computing the hash. \qed
\end{proof}

\begin{intuitionbox}
The two searches give the same big-O for a satisfying reason. An unsuccessful search marches to the end of a length-$\alpha$ list, costing about $\alpha$. A successful search, on average, stops about halfway, costing about $\alpha/2$. Same order, $\Theta(\alpha)$, differing only by the constant $1$ versus $1/2$. And the leading ``$1+$'' in $\Theta(1+\alpha)$ is just the unavoidable cost of computing the hash and touching the array, which we pay even when a list is empty.
\end{intuitionbox}

\begin{corollary}
If $m = \Omega(n)$ -- that is, the number of slots grows at least linearly with the number of keys, so $\alpha = O(1)$ -- then all dictionary operations take $O(1)$ expected time.
\end{corollary}

This is the headline result of the chapter. Keep the load factor bounded by a constant, and a chained hash table does \textsc{Insert}, \textsc{Search}, and \textsc{Delete} in expected constant time. Insert and (pointered) delete are already $O(1)$ in the \emph{worst} case; only search carries the ``expected.''

\begin{complexitybox}
\textbf{Chaining, with $\alpha = n/m$:}
\begin{center}
\begin{tabular}{lll}
\toprule
Operation & Expected & Worst case \\
\midrule
\textsc{Insert} & $O(1)$ & $O(1)$ \\
\textsc{Search} & $\Theta(1+\alpha)$ & $\Theta(n)$ \\
\textsc{Delete} (pointer to $x$) & $O(1)$ & $O(1)$ \\
\bottomrule
\end{tabular}
\end{center}
With $\alpha = O(1)$, search is $O(1)$ expected. Worst-case $\Theta(n)$ happens only if every key collides into one slot -- which good hashing makes astronomically unlikely.
\end{complexitybox}

\begin{labbox}
Build a tiny chained table by hand with $m = 5$ and $h(k) = k \bmod 5$. Insert the keys $12, 7, 33, 20, 18, 25, 2$ in that order, always at the \emph{head} of each list. After all insertions, write out the five lists. (Check, reading each list from its head: slot $0$ should hold $25,20$; slot $2$ should hold $2,7,12$; slot $3$ should hold $18,33$. Slots $1$ and $4$ empty.) Then count how many nodes a successful search for $7$ examines, and an unsuccessful search for $9$.
\end{labbox}

\section{Hash functions: turning keys into slots}

Everything above assumed simple uniform hashing -- an idealized, perfectly uniform spread across slots. Real hash functions are concrete arithmetic, and a bad one clumps keys together and destroys the guarantee. What makes a good hash function?

\begin{focusbox}
A good hash function should (a) be fast to compute, (b) scatter keys so that any clustering or pattern in the keys does \emph{not} translate into clustering in the slots, and (c) ideally satisfy simple uniform hashing. Since we rarely know the exact distribution of keys, we want functions that behave well across many plausible key sets.
\end{focusbox}

Most hash functions first interpret the key as a natural number. A string like \lstinline|"pt"| becomes a number by treating its characters as digits in base $128$: \texttt{p}$=112$, \texttt{t}$=116$, so the key is $112 \cdot 128 + 116 = 14452$. Once a key is a non-negative integer $k$, two classic methods turn it into a slot.

\subsection{The division method}

Pick a table size $m$ and define
\[
  h(k) = k \bmod m.
\]
The hash value is the remainder of $k$ divided by $m$. It is wonderfully fast -- one modulo operation. But the choice of $m$ matters enormously.

\begin{trapbox}
\textbf{Do not let $m$ be a power of $2$ in the division method.} If $m = 2^p$, then $k \bmod m$ is just the low $p$ bits of $k$ -- it ignores all the higher bits entirely. If your keys share patterns in their low bits (common with structured data, pointers, aligned addresses), they collide catastrophically. Powers of $10$ are similarly dangerous for decimal-structured keys.
\end{trapbox}

The folklore fix: choose $m$ to be a \emph{prime} not too close to a power of $2$. A prime modulus mixes all the bits of $k$ into the remainder, so accidental patterns in the keys are far less likely to survive. For example, to store about $2000$ keys with load factor near $3$, you want $m \approx 700$; the prime $701$ is a fine choice.

\begin{example}
Let $m = 701$ and key $k = 1{,}234{,}567$. Then $h(k) = 1{,}234{,}567 \bmod 701$. Compute: $701 \cdot 1761 = 1{,}234{,}461$, and $1{,}234{,}567 - 1{,}234{,}461 = 106$. So the key hashes to slot $106$.
\end{example}

\subsection{The multiplication method}

This method does not care about the exact value of $m$, which frees you to pick $m = 2^p$ for fast indexing. Choose a constant $A$ with $0 < A < 1$. Then
\[
  h(k) = \floor*{\, m \,(k A \bmod 1) \,},
\]
where $kA \bmod 1$ means ``the fractional part of $kA$.'' In words: multiply the key by $A$, throw away the integer part to keep a fraction in $[0,1)$, then scale that fraction up to $[0,m)$ and floor it.

\begin{intuitionbox}
Think of the fractional parts $\{kA\}$ for $k = 1, 2, 3, \dots$ as points dropped on the circle of circumference $1$. If $A$ is a nicely irrational-looking constant, these points spread out evenly around the circle, never bunching up. We then chop the circle into $m$ equal arcs and read off which arc each key landed in. Even, well-mixed arcs mean even, well-mixed slots.
\end{intuitionbox}

Knuth suggested $A \approx (\sqrt{5}-1)/2 = 0.6180339887\ldots$, the reciprocal golden ratio, as a value that spreads keys especially well.

\begin{example}
Let $m = 16$ (so $p=4$), $A = 0.6180339887$, and $k = 100$. Then $kA = 61.80339887$, the fractional part is $0.80339887$, times $m=16$ gives $12.854\ldots$, and the floor is $12$. So $h(100) = 12$.
\end{example}

\begin{marginnotebox}
This material parallels CLRS Chapter 11. Direct addressing and chaining correspond to sections 11.1--11.2; hash functions (division, multiplication) to 11.3; open addressing to 11.4; and universal and perfect hashing to 11.3.3 and 11.5. Knuth's \emph{The Art of Computer Programming}, Volume 3, provides the classic mathematical treatment.
\end{marginnotebox}

\section{Open addressing: no lists at all}

Chaining stores collisions in side lists, which means extra pointers and extra memory beyond the array. \emph{Open addressing} takes a different, leaner stance: \textbf{every element lives in the array itself}, one element per slot, no lists, no external pointers. The load factor $\alpha = n/m$ therefore never exceeds $1$ -- you cannot store more elements than you have slots.

But if a key's natural slot is already occupied by a different key, where does it go? It \emph{probes} a sequence of slots until it finds an empty one. The hash function now takes a second argument, the probe number:
\[
  h : U \times \{0,1,\dots,m-1\} \to \{0,1,\dots,m-1\}.
\]
For a key $k$, we try slots $h(k,0)$, $h(k,1)$, $h(k,2)$, $\dots$ in order. We require that this \emph{probe sequence} $\langle h(k,0), h(k,1), \dots, h(k,m-1)\rangle$ be a permutation of all $m$ slots, so that probing eventually examines every slot and never gives up while empty slots remain.

\begin{algorithm}[H]
\caption{Open-addressing Insert and Search}
\KwIn{Hash table $T[0\,..\,m-1]$, key $k$ (Insert) or sought key $k$ (Search)}
\KwOut{Insert: slot index, or table-full error. Search: slot index or \textsc{nil}.}
\textsc{Hash-Insert}$(T,k)$:\;
\For{$i \gets 0$ \KwTo $m-1$}{
  $j \gets h(k,i)$\;
  \If{$T[j]$ is \textsc{nil} or \textsc{deleted}}{
    $T[j] \gets k$\;
    \Return $j$\;
  }
}
\Return ``table overflow''\;
\BlankLine
\textsc{Hash-Search}$(T,k)$:\;
\For{$i \gets 0$ \KwTo $m-1$}{
  $j \gets h(k,i)$\;
  \lIf{$T[j] = k$}{\Return $j$}
  \lIf{$T[j]$ is \textsc{nil}}{\Return \textsc{nil}}
}
\Return \textsc{nil}\;
\end{algorithm}

Notice the search-stopping rule: we stop and report ``absent'' the moment we hit a genuinely empty (\textsc{nil}) slot, because the key, had it been inserted, would have claimed that slot or an earlier one in its probe sequence. We do \emph{not} stop on an occupied-by-someone-else slot; we keep probing.

\begin{trapbox}
\textbf{Deletion in open addressing is tricky.} You cannot simply set a deleted slot back to \textsc{nil} -- that would create a false ``empty'' in the middle of some other key's probe sequence, causing searches for those keys to stop early and wrongly report ``absent.'' The standard fix is a special \textsc{deleted} marker (a ``tombstone''): search treats it as occupied (keep probing past it), but insert treats it as free (reuse it). The cost: with many tombstones, search times no longer depend only on $\alpha$, which is one reason chaining is often preferred when deletions are frequent.
\end{trapbox}

\subsection{Linear probing}

The simplest probe sequence: if your slot is taken, try the next one, then the next, wrapping around. Given an ordinary hash function $h'(k)$, define
\[
  h(k,i) = \big(h'(k) + i\big) \bmod m.
\]
The probe sequence is $h'(k),\, h'(k)+1,\, h'(k)+2, \dots$ marching forward (mod $m$).

\begin{example}
Let $m = 11$, $h'(k) = k \bmod 11$, and insert $20, 31, 9, 42$ in order.
\begin{itemize}
  \item $20$: $h'=20\bmod 11=9$. Slot $9$ empty $\to$ place at $9$.
  \item $31$: $h'=31\bmod 11=9$. Slot $9$ taken; try $10$ -- empty $\to$ place at $10$.
  \item $9$: $h'=9$. Slot $9$ taken; $10$ taken; try $0$ (wrap) -- empty $\to$ place at $0$.
  \item $42$: $h'=42\bmod 11=9$. Slots $9,10,0$ taken; try $1$ -- empty $\to$ place at $1$.
\end{itemize}
Four keys that all wanted slot $9$ spilled forward into $9,10,0,1$ -- a run of four consecutive filled slots, shown in Figure~\ref{ch11:linear-probing-picture}.
\end{example}

\begin{figure}[h]
\centering
\begin{tikzpicture}[font=\small]
  % index labels and cells for slots 0..10
  \foreach \i/\val in {0/9,1/42,2/{},3/{},4/{},5/{},6/{},7/{},8/{},9/20,10/31} {
    \node[gray] at (0.7*\i,0.55) {\scriptsize \i};
  }
  % filled slots get a tinted box, empty ones a plain box
  \foreach \i/\val in {0/9,1/42,9/20,10/31} {
    \node[draw,fill=trapred!12,minimum width=0.62cm,minimum height=0.55cm] at (0.7*\i,0) {\scriptsize \val};
  }
  \foreach \i in {2,3,4,5,6,7,8} {
    \node[draw,fill=white,minimum width=0.62cm,minimum height=0.55cm] at (0.7*\i,0) {};
  }
  % brace under the run 9,10,0,1 (drawn as two pieces because it wraps)
  \draw[decorate,decoration={brace,amplitude=4pt},trapred] (6.16,-0.45) -- (7.0,-0.45);
  \node[trapred] at (6.6,-0.85) {\scriptsize wrap};
  \draw[decorate,decoration={brace,amplitude=4pt},trapred] (0.7,-0.45) -- (-0.3,-0.45);
  \node[trapred] at (0.2,-0.85) {\scriptsize run of $4$};
\end{tikzpicture}
\caption{Linear probing after inserting $20,31,9,42$ (all with $h'=9$) into $m=11$ slots. They form one contiguous run wrapping past the end: slots $9,10,0,1$.}
\label{ch11:linear-probing-picture}
\end{figure}

That run is the disease called \emph{primary clustering}: long runs of occupied slots build up, and they grow faster than you would guess, because any key hashing anywhere into the run gets appended to its far end, lengthening it. Long runs mean long probes.

\begin{trapbox}
\textbf{Primary clustering.} Linear probing creates long contiguous runs of full slots. A slot just after a long run is very likely to be chosen next (any of the run's keys, plus that slot's own keys, feed it), so runs snowball. Linear probing is cache-friendly -- it walks contiguous memory -- but its clustering makes probe sequences long as $\alpha$ approaches $1$.
\end{trapbox}

\subsection{Quadratic probing}

To break up those runs, jump by a quadratically growing offset instead of $1$:
\[
  h(k,i) = \big(h'(k) + c_1 i + c_2 i^2\big) \bmod m,
\]
for constants $c_1, c_2$ (with $c_2 \neq 0$). Now the probe sequence spreads out: $h', h'+c_1+c_2, h'+2c_1+4c_2, \dots$. Two keys with the \emph{same} initial slot $h'$ still follow the same probe sequence -- a milder defect called \emph{secondary clustering} -- but you avoid the long contiguous runs of linear probing. The constants $c_1, c_2$ and $m$ must be chosen carefully (a common safe choice is $m$ a power of $2$ with $c_1 = c_2 = 1/2$) so that the probe sequence actually visits all $m$ slots.

\subsection{Double hashing}

The best of the simple open-addressing schemes uses a \emph{second} hash function to determine the step size, so that different keys with the same starting slot march off in \emph{different} directions:
\[
  h(k,i) = \big(h_1(k) + i\,h_2(k)\big) \bmod m.
\]
The first probe is at $h_1(k)$; thereafter each step jumps by $h_2(k)$, an amount that \emph{depends on the key}. Two keys sharing a starting slot almost surely have different step sizes, so they diverge immediately -- no secondary clustering. For the probe sequence to hit every slot, $h_2(k)$ must be relatively prime to $m$; a clean way to guarantee this is to make $m$ a power of $2$ and design $h_2$ to always produce odd values, or to make $m$ prime and let $h_2$ return values in $1..m-1$.

\begin{example}
Let $m = 13$, $h_1(k) = k \bmod 13$, and $h_2(k) = 1 + (k \bmod 11)$. Insert key $k = 14$.
$h_1(14) = 1$, $h_2(14) = 1 + (14 \bmod 11) = 1 + 3 = 4$. Probe sequence: $1, 5, 9, 0, 4, 8, 12, 3, 7, 11, 2, 6, 10$ -- step $4$, wrapping mod $13$, visiting all thirteen slots. If slot $1$ is taken, we try $5$; if that is taken, $9$; and so on.
\end{example}

\begin{figure}[h]
\centering
\begin{tikzpicture}[font=\small, scale=0.92, every node/.style={transform shape}]
  % Top coordinate ruler / title
  \node[font=\footnotesize\scshape, text=gray!70!black] at (5.5, 2.5) {Comparison of Open-Addressing Probing Trajectories ($m=11$, initial collision at slot $2$)};

  % Array slots 0 to 10
  \foreach \i in {0,...,10} {
    \ifnum\i=2
      \draw[very thick, draw=trapred, fill=trapred!25, rounded corners=2pt] (\i*1.1, 0) rectangle (\i*1.1+0.85, 0.7);
      \node[trapred!90!black, font=\scriptsize\bfseries] at (\i*1.1+0.425, 0.35) {Coll!};
      \node[trapred!80!black, font=\scriptsize] at (\i*1.1+0.425, -0.3) {$h'(k)=2$};
    \else
      \draw[thick, draw=gray!60, fill=gray!8, rounded corners=2pt] (\i*1.1, 0) rectangle (\i*1.1+0.85, 0.7);
      \node[gray!70!black, font=\scriptsize] at (\i*1.1+0.425, 0.35) {[\i]};
    \fi
  }

  % Linear probing (top, short red hops)
  \draw[-{Latex[length=2.2mm]}, thick, trapred] (2*1.1+0.425, 0.7) to[bend left=30] (3*1.1+0.425, 0.7);
  \draw[-{Latex[length=2.2mm]}, thick, trapred] (3*1.1+0.425, 0.7) to[bend left=30] 
    node[above=1pt, trapred, font=\tiny] {$+1$} (4*1.1+0.425, 0.7);
  \draw[-{Latex[length=2.2mm]}, thick, trapred] (4*1.1+0.425, 0.7) to[bend left=30] 
    node[above=1pt, trapred, font=\tiny] {$+1$} (5*1.1+0.425, 0.7);

  % Quadratic probing (top, high orange hops, clear corridor)
  \draw[-{Latex[length=2.2mm]}, thick, quotientorange, dashed] (2*1.1+0.425, 0.7) to[bend left=60] 
    node[above=2pt, quotientorange, font=\tiny] {$+1^2$} (3*1.1+0.425, 0.7);
  \draw[-{Latex[length=2.2mm]}, thick, quotientorange, dashed] (2*1.1+0.425, 0.7) to[bend left=42] 
    node[above=2pt, quotientorange, font=\tiny] {$+2^2=4$} (6*1.1+0.425, 0.7);

  % Double hashing: key k1 with h2(k1)=3 (bottom, blue leaps)
  \draw[-{Latex[length=2.2mm]}, very thick, sourceblue] (2*1.1+0.425, -0.45) to[bend right=30] 
    node[above=1pt, sourceblue, font=\tiny] {$k_1:\; +3$} (5*1.1+0.425, -0.45);
  \draw[-{Latex[length=2.2mm]}, very thick, sourceblue] (5*1.1+0.425, -0.45) to[bend right=30] 
    node[above=1pt, sourceblue, font=\tiny] {$+3$} (8*1.1+0.425, -0.45);

  % Double hashing: key k2 with h2(k2)=5 (bottom, green leaps)
  \draw[-{Latex[length=2.2mm]}, very thick, targetgreen] (2*1.1+0.425, -0.55) to[bend right=32] 
    node[below=2pt, targetgreen, font=\tiny] {$k_2:\; +5$} (7*1.1+0.425, -0.55);

  % Legend block below
  \node[draw=gray!30, fill=white, rounded corners=3pt, inner sep=5pt, font=\scriptsize, align=left] at (5.5, -3.1) {
    \textbf{Trajectory summary:}\\
    \textcolor{trapred}{\rule{0.6cm}{1.2pt}} \textbf{Linear probing}: $h(k,i) = (2 + i) \bmod 11$ (stride $1$, merges into contiguous runs)\\
    \textcolor{quotientorange}{\textbf{-- --}} \textbf{Quadratic probing}: $h(k,i) = (2 + i^2) \bmod 11$ (accelerating strides, but identical path for any key hashing to $2$)\\
    \textcolor{sourceblue}{\rule{0.6cm}{1.6pt}} \textbf{Double hashing ($k_1$)}: $h(k_1, i) = (2 + 3i) \bmod 11$ \quad
    \textcolor{targetgreen}{\rule{0.6cm}{1.6pt}} \textbf{Double hashing ($k_2$)}: $h(k_2, i) = (2 + 5i) \bmod 11$
  };
\end{tikzpicture}
\caption{Probing trajectories following a collision at slot $h'(k)=2$ in a table of size $m=11$. Linear probing steps through immediate neighbors, causing primary clusters. Quadratic probing accelerates away, but all keys sharing initial slot $2$ trace identical sequences (secondary clustering). Double hashing assigns distinct strides $h_2(k_1)=3$ and $h_2(k_2)=5$, decoupling trajectories across the table.}
\label{ch11:probing-trajectories}
\end{figure}

\begin{intuitionbox}
Linear probing: everybody who collides walks the same single-file line, so jams pile up. Double hashing: each key, even those starting at the same slot, walks at its own stride, so they fan out across the table instead of queuing behind one another. Double hashing uses $\Theta(m^2)$ distinct probe sequences (one per $(h_1,h_2)$ pair), versus only $m$ for linear or quadratic probing, which is why it comes closest to the idealized behavior.
\end{intuitionbox}

\subsection{Cost of open addressing}

The clean analysis assumes \emph{uniform hashing}: each key's probe sequence is equally likely to be any of the $m!$ permutations of the slots. (Double hashing approximates this; linear and quadratic probing do not, but the bound is still a useful guide.) Under this assumption one can prove:

\begin{theorem}[Expected probes in open addressing]
\label{ch11:open-addr-theorem}
In an open-address table with load factor $\alpha = n/m < 1$, under uniform hashing, the expected number of probes is at most $\dfrac{1}{1-\alpha}$ for an unsuccessful search (and for an insertion), and at most $\dfrac{1}{\alpha}\ln\dfrac{1}{1-\alpha}$ for a successful search.
\end{theorem}

\begin{proof}[Intuition for the unsuccessful bound]
A probe sequence of length at least $1$ always happens. It continues to a second probe only if the first slot was occupied, which has probability about $\alpha$ (the fraction of full slots). It reaches a third probe only if the first two were both occupied, probability about $\alpha^2$, and so on. The expected number of probes is therefore about
\[
  1 + \alpha + \alpha^2 + \alpha^3 + \cdots = \frac{1}{1-\alpha},
\]
a geometric series, valid because $\alpha < 1$. \qed
\end{proof}

Put some numbers to that formula. They are the reason hash tables resize before they get too full.

\begin{center}
\begin{tabular}{cccc}
\toprule
$\alpha$ & table is & unsuccessful $\frac{1}{1-\alpha}$ & successful $\frac{1}{\alpha}\ln\frac{1}{1-\alpha}$ \\
\midrule
$0.5$  & half full     & $2.0$  & $1.39$ \\
$0.75$ & 3/4 full      & $4.0$  & $1.85$ \\
$0.9$  & 9/10 full     & $10.0$ & $2.56$ \\
$0.99$ & nearly full   & $100$  & $4.65$ \\
\bottomrule
\end{tabular}
\end{center}

\begin{complexitybox}
\textbf{Open addressing.} Expected probes blow up like $\frac{1}{1-\alpha}$ as $\alpha \to 1$: at $\alpha = 0.9$ an unsuccessful search averages $10$ probes; at $\alpha = 0.99$, a hundred. Practical implementations cap $\alpha$ (often around $0.5$--$0.75$) and \emph{rehash} into a larger table when the cap is hit. Compared to chaining, open addressing uses no pointer overhead and is cache-friendly, but degrades sharply near full and handles deletion awkwardly.
\end{complexitybox}

\begin{walkbox}
Here is something to chew on while you stretch your legs. Under chaining, the expected search cost is $\Theta(1+\alpha)$ and grows \emph{linearly} in $\alpha$ -- doubling the load factor roughly doubles the work. Under open addressing it grows like $\frac{1}{1-\alpha}$, which is fine for small $\alpha$ but \emph{explodes} as $\alpha \to 1$. Why the difference? What is fundamentally happening at $\alpha = 1$ in each scheme? (Hint: in chaining, $\alpha = 1$ just means ``one element per list on average'' -- totally healthy. In open addressing, $\alpha = 1$ means the array is \emph{completely full}, and an unsuccessful search must probe everything.)
\end{walkbox}

\section{Universal hashing: removing the leap of faith}

Everything so far rested on simple uniform hashing -- an \emph{assumption} that our keys happen to spread evenly. But that assumption can fail spectacularly, and not merely through bad luck. If your hash function is fixed and publicly known, an adversary (or simply an unlucky real-world dataset) can hand you $n$ keys that \emph{all} hash to the same slot. Then chaining degenerates to one list of length $n$, and every search is $\Theta(n)$. This is not theoretical paranoia: it has been used to mount real denial-of-service attacks on web servers, by sending requests whose parameters all collide in the server's hash table.

\begin{intuitionbox}
The fix is simple and effective: \emph{do not commit to a single fixed hash function}. At program start, pick a hash function \emph{at random} from a carefully designed family $\mathcal{H}$. Now an adversary, choosing keys, cannot know which function you will use, so they cannot engineer collisions. The randomness lives in \emph{our} choice of $h$, not in the keys -- which means we get a guarantee that holds for \emph{any} input keys, even adversarial ones, in expectation over our coin flips.
\end{intuitionbox}

What property must the family $\mathcal{H}$ have? Exactly this: any two distinct keys should be unlikely to collide, over the random choice of $h$.

\begin{formalbox}
\textbf{Universal family.} A finite collection $\mathcal{H}$ of hash functions, each mapping $U \to \{0,\dots,m-1\}$, is \emph{universal} if for every pair of distinct keys $k \neq \ell$,
\[
  \Pr_{h \in \mathcal{H}}\big[\, h(k) = h(\ell) \,\big] \;\le\; \frac{1}{m},
\]
where the probability is over choosing $h$ uniformly at random from $\mathcal{H}$.
\end{formalbox}

In words: two distinct keys collide with probability at most $1/m$ -- exactly the collision probability you would get from a perfectly random function. And we get that guarantee while still using a small, efficiently-computable family.

\begin{theorem}[Universal hashing gives short chains]
\label{ch11:universal-theorem}
If $h$ is chosen from a universal family $\mathcal{H}$ and we store $n$ keys with chaining in a table of $m$ slots, then for any fixed key $k$, the expected length of the list containing $k$'s slot is at most $1 + \alpha$ if $k$ is in the table, and at most $\alpha$ if it is not. Hence expected search time is $O(1+\alpha)$ -- \emph{for any set of keys}, no assumption required.
\end{theorem}

\begin{proof}
Fix a key $k$ and any other stored key $\ell \neq k$. Define the indicator random variable $X_{k\ell} \in \{0,1\}$ to be $1$ if $h(k) = h(\ell)$ and $0$ otherwise. Because $X_{k\ell}$ takes only values in $\{0,1\}$, its expectation is exactly the collision probability:
\[
  \mathbb{E}[X_{k\ell}] = 1 \cdot \Pr[h(k)=h(\ell)] + 0 \cdot \Pr[h(k) \neq h(\ell)] = \Pr[h(k)=h(\ell)] \le \frac{1}{m},
\]
by the universality of $\mathcal{H}$. The total number of stored keys colliding with $k$ is the sum of these indicators: $Y_k = \sum_{\ell \in K,\, \ell \neq k} X_{k\ell}$. By linearity of expectation -- which holds even though pairwise collision events are not necessarily mutually independent:
\[
  \mathbb{E}[Y_k] = \sum_{\ell \in K,\, \ell \neq k} \mathbb{E}[X_{k\ell}] \le \sum_{\ell \in K,\, \ell \neq k} \frac{1}{m}.
\]
If $k \notin K$, the sum contains exactly $n$ terms, yielding $\mathbb{E}[Y_k] \le n/m = \alpha$. If $k \in K$, the sum excludes $k$ itself and has $n-1$ terms, giving $\mathbb{E}[Y_k] \le (n-1)/m < \alpha$. Adding $1$ for the step of inspecting $k$ itself yields an expected chain traversal of at most $1 + \alpha$. \qed
\end{proof}

This fundamentally strengthens the mathematical guarantee: \textbf{the identical $\Theta(1+\alpha)$ average search cost holds for every possible set of input keys}, because the probability distribution is governed entirely by our own randomized algorithm rather than by assumptions about user inputs.

\subsection{A concrete universal family}

It is not obvious such families exist, so here is one you can implement. Pick a prime $p$ larger than every possible key, so $U \subseteq \{0,1,\dots,p-1\}$. For any $a \in \{1,2,\dots,p-1\}$ and $b \in \{0,1,\dots,p-1\}$, define
\[
  h_{a,b}(k) = \big((a k + b) \bmod p\big) \bmod m.
\]
The family $\mathcal{H}_{p,m} = \{\, h_{a,b} : a \in \{1,\dots,p-1\},\ b \in \{0,\dots,p-1\} \,\}$ -- one function per choice of $(a,b)$ -- is universal. We choose $a$ and $b$ at random when the table is created and keep them fixed thereafter. The proof that collisions happen with probability $\le 1/m$ rests on the fact that, modulo a prime, the map $k \mapsto ak+b$ shuffles the keys around in a way that decouples distinct keys; CLRS section 11.3.3 carries it out in full, and it is a lovely little argument with $\Z_p$ arithmetic.

\begin{example}
Take $p = 17$, $m = 6$, and suppose we randomly drew $a = 3$, $b = 4$. Then $h_{3,4}(8) = ((3\cdot 8 + 4)\bmod 17)\bmod 6 = (28 \bmod 17)\bmod 6 = 11 \bmod 6 = 5$. And $h_{3,4}(15) = ((45+4)\bmod 17)\bmod 6 = (49\bmod 17)\bmod 6 = 15\bmod 6 = 3$. A \emph{different} random draw, say $a=5,b=2$, would send key $8$ to $((42)\bmod 17)\bmod 6 = (8)\bmod 6 = 2$ instead -- different function, different slot. That dependence on the secret $(a,b)$ is precisely what foils an adversary.
\end{example}

\begin{aestheticbox}
Universal hashing is not a textbook curiosity. Around 2003 and again around 2011, researchers showed that fixed, predictable hash functions in popular language runtimes let attackers craft HTTP requests whose keys all collided, turning an $O(1)$ lookup into $O(n)$ and freezing servers under a trickle of traffic -- ``hash-flooding'' denial of service. The defense deployed across Python, Ruby, and others was exactly this chapter's idea: seed the hash function with a per-process random key, so collisions cannot be predicted. The math you just saw is the reason that defense works.
\end{aestheticbox}

\section{Perfect hashing: $O(1)$ worst case for static sets}

Universal hashing gives $O(1)$ \emph{expected} time. Can we ever get $O(1)$ \emph{worst-case} search -- a guarantee with no ``expected'' attached? Yes, in an important special case: when the set of keys is \emph{static} -- fixed and known in advance, never changing afterward. This is the situation for, say, reserved keywords in a programming language, or a fixed dictionary of words. We call the result \emph{perfect hashing}.

\begin{formalbox}
A hash function is \emph{perfect} for a set $K$ of keys if it maps them with \emph{no collisions at all} -- every key gets its own slot. Then a search is a single probe: $O(1)$ worst case, guaranteed.
\end{formalbox}

How do we find a perfect, collision-free function? The clever construction is \emph{two levels of universal hashing}.

\begin{enumerate}
  \item \textbf{Top level.} Use a universal hash function $h$ to spread the $n$ keys into $m = n$ slots. There will still be collisions -- some slots get several keys -- but we don't store them in lists. Instead\dots
  \item \textbf{Second level.} For each slot $j$ that received $n_j$ keys, build a \emph{secondary} hash table of size $m_j = n_j^2$ -- the square of the number of keys in that slot -- using its own universal hash function. Squaring the size is the trick that makes the secondary table collision-free.
\end{enumerate}

\begin{figure}[h]
\centering
\begin{tikzpicture}[font=\small]
  \foreach \i in {0,...,3} {
    \draw (0,-0.7*\i) rectangle (1,-0.7*\i+0.55);
    \node[gray] at (-0.35,-0.7*\i+0.28) {\scriptsize \i};
  }
  \node at (0.5,0.6) {top};
  % slot 0 -> secondary of size 4 (n_0=2)
  \node[draw,fill=targetgreen!12,minimum width=0.5cm] (s0a) at (2.2,0.28) {};
  \node[draw,fill=targetgreen!12,minimum width=0.5cm] (s0b) at (2.7,0.28) {};
  \node[draw,fill=targetgreen!12,minimum width=0.5cm] (s0c) at (3.2,0.28) {};
  \node[draw,fill=targetgreen!12,minimum width=0.5cm] (s0d) at (3.7,0.28) {};
  \draw[-{Latex}] (1,0.28) -- (1.9,0.28);
  \node[gray] at (5.0,0.28) {\scriptsize $n_0{=}2,\ m_0{=}4$};
  % slot 1 -> empty
  \node at (2.0,-0.42) {\textsc{nil}};
  % slot 2 -> secondary size 1 (n_2=1)
  \node[draw,fill=quotientorange!15,minimum width=0.5cm] (s2) at (2.2,-1.12) {};
  \draw[-{Latex}] (1,-1.12) -- (1.9,-1.12);
  \node[gray] at (4.0,-1.12) {\scriptsize $n_2{=}1,\ m_2{=}1$};
  % slot 3 -> secondary size 9 (n_3=3)
  \foreach \k in {0,...,8} {
    \node[draw,fill=sourceblue!12,minimum width=0.32cm] at (2.0+0.4*\k,-1.82) {};
  }
  \draw[-{Latex}] (1,-1.82) -- (1.9,-1.82);
  \node[gray] at (6.6,-1.82) {\scriptsize $n_3{=}3,\ m_3{=}9$};
\end{tikzpicture}
\caption{Perfect hashing: a top-level table whose slots point to small secondary tables of size $n_j^2$, each collision-free.}
\label{ch11:perfect-hashing-picture}
\end{figure}

Two questions should be bothering you. First, why does squaring make a slot collision-free? And second, isn't all that squared space enormous?

\begin{lemma}[Squaring kills collisions, in expectation]
\label{ch11:perfect-lemma}
If you hash $n_j$ keys into a table of size $m_j = n_j^2$ using a universal hash function, the expected number of colliding pairs is less than $1/2$. So a single random choice of function is collision-free with probability greater than $1/2$, and after a few tries you are guaranteed to find one.
\end{lemma}

\begin{proof}[Why]
There are $\binom{n_j}{2} < n_j^2/2$ pairs of keys. Each pair collides with probability at most $1/m_j = 1/n_j^2$ by universality. By linearity of expectation, the expected number of colliding pairs is less than $\frac{n_j^2}{2}\cdot\frac{1}{n_j^2} = \frac12$. A non-negative quantity with expectation below $1/2$ must be $0$ more than half the time, so a random universal function gives zero collisions with probability $> 1/2$. Keep drawing until it does. \qed
\end{proof}

That answers the first question: each secondary table, sized as the square of its occupancy, is collision-free, so a search there is a single probe. A perfect-hash search is thus: one probe at the top level to find the right secondary table, one probe inside it -- two probes total, $O(1)$ \emph{worst case}.

Now the space. Surely $\sum_j n_j^2$ could be huge -- if all $n$ keys piled into one slot, it would be $n^2$. The saving grace is that, by choosing the top-level function from a universal family, that pile-up almost never happens.

\begin{theorem}[Perfect hashing uses linear space]
\label{ch11:perfect-space-theorem}
If we choose the top-level universal hash function with $m = n$, then the expected total size of all secondary tables is less than $2n$. Hence perfect hashing stores $n$ keys in $O(n)$ space while guaranteeing $O(1)$ worst-case search.
\end{theorem}

The key fact behind this is that $\mathbb{E}\big[\sum_j n_j^2\big] < 2n$ when $m = n$ -- the expected sum of squared bucket sizes is linear, not quadratic, precisely because universality keeps the buckets balanced. (CLRS section 11.5 proves it; it again comes down to counting colliding pairs.) As with the secondary tables, we simply redraw the top-level function if the total turns out too large; on average a couple of tries suffice.

\begin{focusbox}
\textbf{Perfect hashing in brief:} use two levels of universal hashing. The top level spreads $n$ keys over $n$ slots; each slot's $n_j$ keys get a private secondary table of size $n_j^2$, which universality makes collision-free. Result: $O(1)$ \emph{worst-case} lookup using $O(n)$ space -- but only for a \emph{static} key set, because rebuilding on every insert would be too slow.
\end{focusbox}

\begin{historybox}
The dictionary problem arose in the earliest commercial computers at IBM. In 1953, Hans Peter Luhn was tasked with indexing chemical formulas and bibliographic records on punch cards. Direct addressing was impossible: the universe of chemical compounds was astronomical, yet RAM was measured in tiny magnetic core words. Luhn proposed computing a numerical address from the key by folding and division, inventing hash addressing. Simultaneously, A.~D.~Linh (1953) developed separate chaining with linked lists to resolve collisions, and Arnold I.~Dumey (1956) formulated open addressing with linear probing. Two decades later, J.~Lawrence Carter and Mark N.~Wegman (1979) introduced universal hashing, providing the first mathematically rigorous defense against worst-case collision attacks. Two-level perfect hashing was achieved by Michael Fredman, J\'anos Koml\'os, and Endre Szemer\'edi in 1984 (the FKS scheme).

Early practitioners treated hash functions as ad-hoc engineering hacks, assuming real-world inputs would distribute uniformly like independent random variables. This naive assumption repeatedly failed: systematic patterns in input data (such as zero-padded customer IDs or common compiler identifier prefixes) produced severe clustering. In open addressing, linear probing suffered from primary clustering, where contiguous blocks of filled slots acted as magnets for subsequent keys. Worse, fixed public hash functions created severe denial-of-service vulnerabilities: web servers running PHP, Java, and Python could be brought to a halt by sending HTTP requests with thousands of distinct parameter names crafted to collide into a single hash bucket, degenerating $O(1)$ lookups into $\Theta(n)$ linked-list traversals. Universal hashing solved this by drawing a randomized hash function fresh at runtime from a provably collision-bounded family.

Expected $O(1)$ dictionary search, insertion, and deletion revolutionized software architecture. Hash tables power compiler symbol tables, database indices, associative arrays in modern scripting languages (Python dicts, JavaScript objects), deduplication caches, and distributed key-value stores. Yet this carries an honest trade-off: hash tables completely destroy key ordering: they cannot support \textsc{Minimum}, \textsc{Maximum}, \textsc{Predecessor}, \textsc{Successor}, or sorted range queries ($[\text{low}, \text{high}]$) without scanning all $m$ slots ($\Theta(m)$ time). Dynamic resizing causes latency spikes when tables double, and linked-list chaining incurs cache misses from pointer chasing. To restore order without sacrificing dynamic efficiency, we turn to tree-based structures. In \textbf{Chapter~12 (Binary Search Trees)}, we abandon pseudo-random scattering in favor of the binary-search-tree invariant, unlocking monotonic in-order traversals and logarithmic predecessor searches.
\end{historybox}

\section{Probabilistic Data Structures: Bloom Filters and Count-Min Sketch}
\label{ch11:probabilistic}

Every hash table we have analyzed up to this point shares an implicit contract: it stores the keys themselves, or pointers to where those keys live. If an item is present, the table returns its record; if it is absent, the table reports absence with absolute certainty.

Now consider what happens when that contract collides with large-scale systems engineering.

Imagine an LSM-tree (Log-Structured Merge-tree) powering a distributed storage engine like Google Bigtable, Apache Cassandra, or RocksDB. The system stores billions of records across hundreds of immutable files on disk, known as SSTables (Sorted String Tables). When a client requests a read for key $k$, the database might need to consult forty different SSTables. Must it perform forty separate disk reads just to discover that key $k$ is absent from thirty-nine of them? A single solid-state drive or spinning disk read takes anywhere from tens of microseconds to several milliseconds. Multiplying that latency by forty disk accesses per read would cripple query throughput.

A natural first reaction is: keep an in-memory hash table of all keys contained in each SSTable. But compute the arithmetic. If a storage node hosts $10^9$ keys and each key is a 64-bit integer, storing the raw keys alone requires:
\[
  10^9 \times 8\text{ bytes} = 8\text{ gigabytes of RAM}.
\]
Once you include bucket arrays, collision pointers, and load-factor padding to keep lookups fast, the memory footprint easily balloons to $16$ to $32$ gigabytes. For an embedded database, a network router line-card, or a web browser validating visited URLs against a blacklist of hundreds of millions of malicious phishing domains, keeping all keys in RAM is an impossible luxury.

The breakthrough is to ask an uncomfortable question: do we really need $100\%$ certainty in both directions?

\subsection{Bloom filters: trading perfection for compactness}

In 1970, Burton Howard Bloom answered with an elegant structural compromise. Suppose we do not demand complete certainty when an item is reported present, but we maintain an ironclad guarantee when an item is reported absent:
\begin{itemize}
  \item If the structure reports \textbf{``Definitely Absent''}, the item is guaranteed not to be in the set. We can skip the disk seek entirely, with zero risk of missing real data.
  \item If the structure reports \textbf{``Probably Present''}, the item might be in the set, but there is a tiny, tunable probability $\eps$ (such as $1\%$) that it is a \emph{false positive}.
\end{itemize}
This asymmetric contract provides an absolute guarantee of \textbf{zero false negatives} ($\P(\text{report absent} \mid \text{present}) = 0$) in exchange for a small, controlled rate of \textbf{false positives} ($\P(\text{report present} \mid \text{absent}) = \eps$). If $\eps = 0.01$, ninety-nine out of every one hundred queries for non-existent keys are eliminated in RAM without touching disk. The one remaining query incurs a disk read that returns empty, a harmless performance tax.

How can a data structure provide this guarantee without storing any keys? Instead of storing the keys, a \emph{Bloom filter} stores only a single flat array $B[0\,..\,m-1]$ of $m$ bits, initially all set to $0$, accompanied by $k$ independent, uniform hash functions:
\[
  h_1, h_2, \dots, h_k : U \to \{0, 1, \dots, m-1\}.
\]
The mechanics are shown in Figure~\ref{ch11:fig-bloom-filter}(a):
\begin{itemize}
  \item \textbf{Insertion of key $x$:} Compute the $k$ hash values $h_1(x), h_2(x), \dots, h_k(x)$ and set every corresponding bit to $1$:
  \[
    B[h_i(x)] \gets 1 \quad \text{for each } i \in \{1, 2, \dots, k\}.
  \]
  \item \textbf{Query for key $y$:} Inspect the $k$ probe positions $h_1(y), h_2(y), \dots, h_k(y)$ in array $B$:
  \begin{itemize}
    \item If \emph{at least one} bit $B[h_i(y)] = 0$, then $y$ was never inserted. Had $y$ been inserted, that bit would have been flipped to $1$. The query immediately short-circuits and reports \textsc{False} (``Definitely Absent'').
    \item If \emph{all} $k$ bits $B[h_i(y)] = 1$, the query reports \textsc{True} (``Probably Present'').
  \end{itemize}
\end{itemize}
Deletion is impossible in a standard Bloom filter: clearing a bit back to $0$ could erase the presence of other keys whose hash functions happened to land on that same slot.

\subsection{First-principles derivation of false positives and optimal parameters}

Let us derive the mathematics of the false-positive probability $\eps$ directly from first principles.

Suppose we insert $n$ distinct keys into an array of $m$ bits using $k$ independent hash functions that scatter each key uniformly across the $m$ slots.

Consider a single fixed bit in the array during one hash computation of one key. The probability that this specific bit is \emph{not} chosen by the hash function is:
\[
  1 - \frac{1}{m}.
\]
Because each of the $n$ keys is hashed by $k$ independent hash functions, there are $kn$ total bit selections made during the insertion phase. Assuming independence across hash choices, the probability that our chosen bit remains $0$ after all $n$ insertions have taken place is:
\[
  p_0 = \left(1 - \frac{1}{m}\right)^{kn} = \left[\left(1 - \frac{1}{m}\right)^m\right]^{kn/m}.
\]
Using the fundamental limit $\lim_{m \to \infty} (1 - 1/m)^m = e^{-1}$, we obtain the clean approximation:
\[
  p_0 \approx e^{-kn/m}.
\]
The probability that this bit has been flipped to $1$ is therefore:
\[
  p_1 = 1 - p_0 \approx 1 - e^{-kn/m}.
\]
Now consider querying a key $y$ that was never inserted into the set. What is the probability that the filter mistakenly answers that $y$ is present? A false positive occurs if and only if all $k$ bits examined during the query of $y$ happen to be $1$. Assuming the $k$ probe locations behave independently, the false-positive probability $\eps$ is:
\begin{equation}
  \label{ch11:eq-bloom-fp}
  \eps = (p_1)^k \approx \left(1 - e^{-kn/m}\right)^k.
\end{equation}
Look closely at equation~\eqref{ch11:eq-bloom-fp}. For a fixed ratio of bits per key, $m/n$, what value of $k$ minimizes the error rate $\eps$?
\begin{itemize}
  \item If $k$ is very small (say $k=1$), we test only a single bit. An absent key easily lands on a bit that was set by some other key, yielding a high false-positive rate.
  \item If $k$ is very large, each inserted key sets many bits. The array quickly fills with ones, so $p_1 \to 1$, and again almost every query returns a false positive.
\end{itemize}
Between these two extremes lies a unique minimum, visible in Figure~\ref{ch11:fig-bloom-filter}(b). To calculate this minimum, take the natural logarithm of $\eps$:
\[
  g(k) = \ln \eps = k \ln\left(1 - e^{-kn/m}\right).
\]
Let $p = e^{-kn/m}$, the fraction of bits that remain $0$. Taking the natural logarithm gives $\ln p = -kn/m$, so:
\[
  k = -\frac{m}{n} \ln p.
\]
Substituting this expression for $k$ into $g(k)$:
\[
  g(k) = -\frac{m}{n} \ln p \cdot \ln(1 - p).
\]
Because $-m/n$ is a negative constant, minimizing $g(k)$ is equivalent to maximizing the product:
\[
  f(p) = \ln p \cdot \ln(1 - p) \quad \text{for } p \in (0, 1).
\]
Notice the symmetry of $f(p)$: swapping $p$ with $1-p$ leaves the product unchanged ($f(p) = f(1-p)$). The unique maximum on the open interval $(0, 1)$ occurs at the midpoint:
\[
  p = \frac{1}{2}.
\]
This yields a memorable physical conclusion: \textbf{a Bloom filter operates at its theoretical optimum when exactly half the bits are zeros and half the bits are ones!}

Setting $p = e^{-kn/m} = 1/2$, we solve for the optimal number of hash functions $k^\star$:
\[
  -\frac{k n}{m} = \ln\left(\frac{1}{2}\right) = -\ln 2 \implies k^\star = \frac{m}{n} \ln 2 \approx 0.6931 \left(\frac{m}{n}\right).
\]
When operating at $k^\star$, the minimized false-positive probability is:
\[
  \eps = \left(1 - \frac{1}{2}\right)^{k^\star} = 2^{-k^\star} = 2^{-(m/n)\ln 2} = e^{-(m/n)(\ln 2)^2} \approx (0.6185)^{m/n}.
\]
Rearranging this equation gives the required memory ratio $m/n$ as a function of the desired error rate $\eps$:
\[
  \frac{m}{n} = -\frac{\ln \eps}{(\ln 2)^2} \approx -1.4427 \ln \eps \approx 2.081 \log_2\left(\frac{1}{\eps}\right).
\]
Consider what this means in practice:
\begin{itemize}
  \item For a $1\%$ false-positive rate ($\eps = 0.01$), the formula requires $m/n \approx 1.4427 \times \ln(100) \approx 9.57 \approx 10$ bits per key, with $k^\star = 10 \ln 2 \approx 7$ hash functions.
  \item For a $0.1\%$ false-positive rate ($\eps = 0.001$), we need $m/n \approx 1.4427 \times \ln(1000) \approx 14.35 \approx 15$ bits per key, with $k^\star = 15 \ln 2 \approx 10$ hash functions.
\end{itemize}
At $10$ bits per key (just $1.25$ bytes), a billion keys require only $1.2$ gigabytes of memory, fitting comfortably in RAM on standard servers while filtering out $99\%$ of unnecessary disk reads.

\begin{figure}[htbp]
\centering
\begin{tikzpicture}[font=\small, >=Latex]
  % -------------------------------------------------------------
  % Panel (a): 12-bit Bloom filter insertion and membership query
  % -------------------------------------------------------------
  \begin{scope}[xshift=0cm, yshift=0cm]
    \node[anchor=west, font=\small\bfseries] at (-0.4, 3.55) {(a) 12-bit Bloom filter ($k=3$)};

    % Key insertion sources at the top
    \node[draw=sourceblue, fill=sourceblue!12, rounded corners=3pt, inner sep=3pt, font=\footnotesize\bfseries] (keyx) at (1.2, 2.4) {Insert $x$};
    \node[draw=targetgreen!80!black, fill=targetgreen!15, rounded corners=3pt, inner sep=3pt, font=\footnotesize\bfseries] (keyy) at (4.2, 2.4) {Insert $y$};

    % Bit array cells (12 bits: 0 to 11) with value and index inside
    \def\cw{0.46}
    \foreach \idx/\val/\col in {
      0/0/white,
      1/1/sourceblue!25,
      2/0/white,
      3/0/white,
      4/1/structurepurple!30,
      5/0/white,
      6/0/white,
      7/1/targetgreen!25,
      8/0/white,
      9/1/sourceblue!25,
      10/0/white,
      11/1/targetgreen!25
    } {
      \draw[thick, fill=\col] (\idx*\cw, 0.65) rectangle ++(\cw, 0.6);
      \node[font=\footnotesize\bfseries] at (\idx*\cw + 0.5*\cw, 1.02) {\val};
      \node[font=\tiny, text=gray!80!black] at (\idx*\cw + 0.5*\cw, 0.8) {\idx};
    }
    \node[anchor=east, font=\scriptsize\bfseries] at (-0.05, 1.02) {$B$:};
    \node[anchor=east, font=\tiny, text=gray!80!black] at (-0.05, 0.8) {idx:};

    % Arrows from x to slots 1, 4, 9
    \draw[->, thick, sourceblue] (keyx.south) to[out=270, in=90] (1*\cw + 0.5*\cw, 1.25);
    \draw[->, thick, sourceblue] (keyx.south) to[out=270, in=90] (4*\cw + 0.3*\cw, 1.25);
    \draw[->, thick, sourceblue] (keyx.south) to[out=270, in=90] (9*\cw + 0.5*\cw, 1.25);

    % Arrows from y to slots 4, 7, 11
    \draw[->, thick, targetgreen!80!black] (keyy.south) to[out=270, in=90] (4*\cw + 0.7*\cw, 1.25);
    \draw[->, thick, targetgreen!80!black] (keyy.south) to[out=270, in=90] (7*\cw + 0.5*\cw, 1.25);
    \draw[->, thick, targetgreen!80!black] (keyy.south) to[out=270, in=90] (11*\cw + 0.5*\cw, 1.25);

    % Queries at the bottom
    \node[draw=trapred, fill=trapred!10, rounded corners=3pt, inner sep=3pt, font=\scriptsize] (q1) at (0.8, -0.85) {Query $z_1$};
    \node[draw=quotientorange!90!black, fill=quotientorange!15, rounded corners=3pt, inner sep=3pt, font=\scriptsize] (q2) at (4.6, -0.85) {Query $z_2$};

    % Arrows from queries up to probed slots (landing cleanly at bottom of cell 0.65)
    \draw[->, thick, dashed, trapred] (q1.north) to[out=90, in=270] (2*\cw + 0.5*\cw, 0.65);
    \draw[->, thick, dashed, trapred!60] (q1.north) to[out=90, in=270] (7*\cw + 0.2*\cw, 0.65);
    \draw[->, thick, dashed, trapred!60] (q1.north) to[out=90, in=270] (9*\cw + 0.2*\cw, 0.65);

    \draw[->, thick, dashed, quotientorange!90!black] (q2.north) to[out=90, in=270] (1*\cw + 0.8*\cw, 0.65);
    \draw[->, thick, dashed, quotientorange!90!black] (q2.north) to[out=90, in=270] (7*\cw + 0.8*\cw, 0.65);
    \draw[->, thick, dashed, quotientorange!90!black] (q2.north) to[out=90, in=270] (11*\cw + 0.2*\cw, 0.65);

    % Explanatory result callouts below
    \node[anchor=north, align=center, font=\tiny, text=trapred!90!black] at (0.8, -1.25) {
      $B[2]=0 \implies$ \textbf{Absent}\\
      ($0$ false negatives)
    };
    \node[anchor=north, align=center, font=\tiny, text=quotientorange!90!black] at (4.6, -1.25) {
      $B[1,7,11]=1 \implies$ \textbf{False Positive}\\
      (Never inserted)
    };
  \end{scope}

  % -------------------------------------------------------------
  % Panel (b): False-positive probability curve as function of k
  % -------------------------------------------------------------
  \begin{scope}[xshift=7.6cm, yshift=0cm]
    \node[anchor=west, font=\small\bfseries] at (-0.4, 3.55) {(b) False-positive curve $\eps(k) = (1 - e^{-kn/m})^k$};

    % Coordinate axes
    \draw[->, thick] (0, 0) -- (5.0, 0) node[right, font=\footnotesize] {$k$};
    \draw[->, thick] (0, 0) -- (0, 3.0) node[left=2pt, font=\footnotesize] {$\eps$};

    % Axis ticks
    \foreach \k/\klabel in {1/1, 2/2, 4/4, 6/6, 8/8} {
      \draw (\k*0.52, 0.06) -- (\k*0.52, -0.06) node[below, font=\scriptsize] {\klabel};
    }
    \foreach \y/\ylabel in {0.7/0.05, 1.4/0.10, 2.1/0.15, 2.8/0.20} {
      \draw (0.06, \y) -- (-0.06, \y) node[left, font=\scriptsize] {\ylabel};
    }

    % Theoretical curve for m/n = 8 (k* = 8*ln 2 approx 5.54 -> x = 5.54*0.52 = 2.88)
    \draw[very thick, sourceblue]
      (0.52, 2.7) % k=1
      .. controls (1.1, 0.95) and (1.9, 0.35) .. (2.88, 0.28) % minimum at k*
      .. controls (3.6, 0.32) and (4.2, 0.65) .. (4.6, 1.45); % k=8+

    % Minimum point k*
    \fill[trapred] (2.88, 0.28) circle (2pt);
    \draw[dashed, trapred, thick] (2.88, 0) -- (2.88, 0.28);
    \node[below, font=\scriptsize\bfseries, text=trapred] at (2.88, -0.38) {$k^\star = \frac{m}{n}\ln 2$};

    % Annotations placed safely above curve
    \node[anchor=south west, font=\tiny, align=left, text=sourceblue!80!black] at (1.1, 1.45) {
      \textbf{Too few hashes:}\\
      Under-probed
    };

    \node[anchor=south west, font=\tiny, align=left, text=sourceblue!80!black] at (3.3, 1.45) {
      \textbf{Too many hashes:}\\
      Array $> 50\%$ ones
    };

    \node[anchor=north west, font=\tiny, align=left, draw=targetgreen!70!black, fill=targetgreen!8, rounded corners=2pt, inner sep=2.5pt] at (0.8, 2.95) {
      At $k^\star$: array is \textbf{50\% ones}\\
      $\eps_{\min} = (1/2)^{k^\star} \approx (0.6185)^{m/n}$
    };
  \end{scope}
\end{tikzpicture}
\caption{The mechanics and mathematics of a Bloom filter. (a) A 12-bit array populated by keys $x$ and $y$ with $k=3$. Querying $z_1$ immediately reveals it is absent because $B[2] = 0$, guaranteeing zero false negatives. Querying $z_2$ yields a false positive because its three probe positions ($1, 7, 11$) happen to have been set by other elements. (b) The false-positive probability $\eps$ as a function of the number of hash functions $k$. The curve attains its global minimum at $k^\star = (m/n)\ln 2$, the operating point where exactly half the array bits are ones.}
\label{ch11:fig-bloom-filter}
\end{figure}

\subsection{From set membership to stream frequency: Count-Min Sketch}

A Bloom filter answers a binary question about set membership: has key $x$ appeared at least once?

In high-throughput distributed systems -- such as network switches inspecting packet headers, search engines tracking surging queries, and telemetry pipelines aggregating metrics over millions of events per second -- we face a quantitative question:
\[
  \text{``In a streaming sequence of billions of items, how many times has item } x \text{ appeared?''}
\]
Storing an exact counter for every unique item in a hash table requires memory proportional to the number of distinct keys. When the universe of keys contains millions of IP addresses or URLs, an exact table exceeds cache limits.

In 2005, Graham Cormode and S.~Muthukrishnan introduced the \emph{Count-Min Sketch}, extending the bit-array philosophy of the Bloom filter to integer frequency estimation.

Instead of a single one-dimensional bit array, a Count-Min Sketch maintains a two-dimensional grid of integer counters $C[1\,..\,d, 1\,..\,w]$ with depth $d$ and width $w$:
\[
  \begin{pmatrix}
    C[1, 1] & C[1, 2] & \dots & C[1, w] \\
    C[2, 1] & C[2, 2] & \dots & C[2, w] \\
    \vdots  & \vdots  & \ddots & \vdots  \\
    C[d, 1] & C[d, 2] & \dots & C[d, w]
  \end{pmatrix}
\]
Each row $i \in \{1, \dots, d\}$ is paired with its own independent hash function $h_i : U \to \{1, \dots, w\}$. All counters are initialized to zero.

\begin{itemize}
  \item \textbf{Update on arrival of item $x$:} Whenever an occurrence of item $x$ arrives in the stream (or its count increases by $c \ge 1$), we increment the hashed counter in every row:
  \[
    C[i, h_i(x)] \gets C[i, h_i(x)] + c \quad \text{for each row } i \in \{1, \dots, d\}.
  \]
  \item \textbf{Point query for estimated frequency $\hat{f}(x)$:} To estimate the true frequency $f(x)$ of item $x$, inspect the counter at column $h_i(x)$ across all $d$ rows, and return the \textbf{minimum}:
  \begin{equation}
    \label{ch11:eq-count-min}
    \hat{f}(x) = \min_{1 \le i \le d} C[i, h_i(x)].
  \end{equation}
\end{itemize}

Why does taking the minimum work so well?

Because counters are only ever incremented by non-negative amounts, hash collisions can only ever add positive excess to a counter. If another item $y \neq x$ collides with $x$ in row $i$ ($h_i(y) = h_i(x)$), $y$'s occurrences inflate $C[i, h_i(x)]$. Collisions never subtract. Consequently, every single cell is an upper bound on the true count:
\[
  C[i, h_i(x)] \ge f(x) \quad \text{for every row } i \in \{1, \dots, d\}.
\]
Because every counter satisfies $C[i, h_i(x)] \ge f(x)$, their minimum also satisfies $\hat{f}(x) \ge f(x)$. The estimate never underestimates. By taking the minimum over $d$ rows hashed independently, we pick the row whose counter suffered the least collision contamination from other elements.

\begin{theorem}[Count-Min Sketch accuracy guarantee]
Let $N = \sum_{y} f(y)$ be the total number of items processed in the stream. If we choose width $w = \lceil e / \eps \rceil$ and depth $d = \lceil \ln(1 / \delta) \rceil$, the estimate $\hat{f}(x)$ satisfies:
\[
  f(x) \le \hat{f}(x) \le f(x) + \eps N
\]
with probability at least $1 - \delta$.
\end{theorem}

\begin{proof}
Fix an item $x$ and consider row $i$. The counter $C[i, h_i(x)]$ equals the true frequency $f(x)$ plus the contributions from all colliding items $y \neq x$:
\[
  C[i, h_i(x)] = f(x) + \sum_{y \neq x} f(y) \cdot \mathbb{I}\{h_i(y) = h_i(x)\}.
\]
Define the random indicator variable $X_{i, y} = \mathbb{I}\{h_i(y) = h_i(x)\}$. Under uniform hashing into $w$ slots, $\mathbb{E}[X_{i, y}] = \P(h_i(y) = h_i(x)) = 1/w$. By linearity of expectation, the expected overcount in row $i$ is:
\[
  \mathbb{E}\big[C[i, h_i(x)] - f(x)\big] = \sum_{y \neq x} f(y) \cdot \frac{1}{w} \le \frac{N}{w}.
\]
With $w = \lceil e / \eps \rceil$, this expected error is at most $\eps N / e$. Because the overcount $C[i, h_i(x)] - f(x)$ is non-negative, Markov's inequality implies:
\[
  \P\Big(C[i, h_i(x)] - f(x) \ge \eps N\Big) \le \frac{\mathbb{E}[C[i, h_i(x)] - f(x)]}{\eps N} \le \frac{\eps N / e}{\eps N} = \frac{1}{e}.
\]
For the minimum $\hat{f}(x)$ to exceed $f(x) + \eps N$, \emph{every single row} $i \in \{1, \dots, d\}$ must independently fail to stay within the error bound. Because the $d$ rows use independent hash functions:
\[
  \P\Big(\hat{f}(x) > f(x) + \eps N\Big) \le \left(\frac{1}{e}\right)^d = e^{-d}.
\]
Setting $d = \lceil \ln(1 / \delta) \rceil$ yields $e^{-d} \le \delta$. Therefore, with probability at least $1 - \delta$, the estimate is bounded by $f(x) \le \hat{f}(x) \le f(x) + \eps N$. \qed
\end{proof}

The defining property of both Bloom filters and Count-Min sketches is that their
memory footprint depends only on the desired accuracy parameters ($\eps$ and $\delta$),
independent of the length of the data stream or the size of the key universe. By
accepting a small, bounded probability of error, they track immense streams in small,
predictable memory.

\section{Choosing a scheme: a practical recap}

You now have several collision strategies. Which to reach for?

\begin{itemize}
  \item \textbf{Chaining} is the safe default. It tolerates load factors above $1$, degrades gracefully, and handles deletions cleanly. Its cost is pointer overhead and worse cache behavior. Pair it with a \emph{universal} hash function and you get $O(1)$ expected time with no assumptions on the keys -- and protection against hash-flooding.
  \item \textbf{Open addressing} (especially \emph{double hashing}) is leaner and cache-friendlier, with no pointer overhead. Use it when memory is tight and deletions are rare, and keep the load factor well under $1$ (resize at, say, $0.7$). Watch out for tombstones and clustering.
  \item \textbf{Perfect hashing} is the specialist: a \emph{static}, known-in-advance key set where you want a guaranteed single-probe lookup -- compiler keyword tables, immutable lookup sets.
  \item \textbf{Bloom filters and sketches} are the probabilistic specialists. When datasets contain billions of keys that cannot physically fit in RAM, or when monitoring high-velocity data streams, trade an infinitesimal probability of false positives for zero false negatives and dramatic memory savings ($\approx 10$ bits per key for $1\%$ error with a Bloom filter, and counter grids for frequency estimation with Count-Min sketches).
\end{itemize}

\begin{recoverybox}
If the chapter feels like a lot, here is the spine to hold onto. (1) Direct addressing is $O(1)$ but needs $|U|$ space -- too much. (2) Hashing shrinks $|U|$ down to $m$ slots, creating collisions. (3) Chaining resolves collisions with lists; expected search is $\Theta(1+\alpha)$. (4) Open addressing stores everything in the array and probes; expected search is about $\frac{1}{1-\alpha}$. (5) Universal hashing makes the expectation hold for \emph{any} keys by randomizing the function. (6) Perfect hashing achieves $O(1)$ worst-case for static sets via two universal levels. (7) Bloom filters trade a tunable false-positive rate for zero false negatives, compressing set membership into $\approx 10$ bits per key; Count-Min sketches generalize this to estimate stream frequencies using counter grids. Everything else is detail and proof.
\end{recoverybox}

Notice the ultimate trade-off hash tables make in exchange for expected $O(1)$ speed:
they destroy key order. By scattering keys uniformly across an array, a hash table
cannot answer questions like ``what is the minimum key?'', ``which key comes next
after $x$?'', or ``list all elements between $10$ and $20$'' without scanning all $m$
slots. When your workload requires dynamic insertions and deletions while keeping keys
in sorted order, you need a different structure. In Chapter~12, we turn to
\textbf{Binary Search Trees}, where hierarchical branching preserves key order to
support searches, successors, and sorted range queries in logarithmic time.

\section{Exercises}

\begin{exercise}
A direct-address table $T[0\,..\,9]$ starts empty. Show its contents (which slots are non-\textsc{nil}) after \textsc{Insert} of keys $3$, $7$, $1$, $7$ in order, then \textsc{Delete} of key $7$. (Treat re-inserting $7$ as overwriting its slot.) What is the worst-case time of each operation?
\end{exercise}

\begin{exercise}
A chained hash table has $m = 7$ slots and uses $h(k) = k \bmod 7$. Insert keys $15, 22, 8, 1, 29, 36$ in order, always at the head of the list. Draw all seven lists. What is the load factor? For a successful search of $22$, how many list nodes are examined? For an unsuccessful search of $9$?
\end{exercise}

\begin{exercise}
You expect to store about $n = 10{,}000$ keys and want an average chain length of about $2$. What table size $m$ should you choose, and why should it be a prime not near a power of $2$ if you use the division method? Give a concrete suitable prime.
\end{exercise}

\begin{exercise}
Using the multiplication method with $A = (\sqrt 5 - 1)/2 \approx 0.6180339887$ and $m = 8$, compute $h(k)$ for $k = 50$ and $k = 51$. Show the intermediate fractional part each time. Do consecutive keys land in nearby slots or scattered ones?
\end{exercise}

\begin{exercise}
A table of size $m = 11$ uses linear probing with $h'(k) = k \bmod 11$. Insert the keys $26, 37, 15, 48, 4$ in order and show the final array. Identify the longest run of consecutive occupied slots, and explain in your own words why such runs form (primary clustering).
\end{exercise}

\begin{exercise}
With $m = 13$, $h_1(k) = k \bmod 13$, and $h_2(k) = 1 + (k \bmod 11)$, write out the first five probe positions for double hashing of key $k = 27$. Verify that $h_2(27)$ is between $1$ and $11$ and therefore coprime to $13$.
\end{exercise}

\begin{exercise}
Using Theorem~\ref{ch11:open-addr-theorem}, compute the expected number of probes for an \emph{unsuccessful} search when $\alpha = 0.8$ and when $\alpha = 0.95$. Then explain, in terms of these numbers, why production hash tables resize once the load factor crosses a threshold like $0.7$.
\end{exercise}

\begin{exercise}
Consider the universal family $h_{a,b}(k) = ((ak+b)\bmod p)\bmod m$ with $p = 17$ and $m = 6$. With $a = 2$, $b = 3$, compute $h_{2,3}(k)$ for $k = 0, 5, 11$. Then pick a \emph{different} pair $(a,b)$ and recompute $h(5)$; observe that the slot changes. Why does the randomness in $(a,b)$ defend against an adversary who chooses the keys?
\end{exercise}

\begin{exercise}
Prove that if a hash function $h$ maps a set of $n$ keys into $m$ slots and $n > m$, then there must be at least one collision. (This is the pigeonhole principle. State it carefully and connect it to why collision resolution is unavoidable when $|K| > m$.)
\end{exercise}

\begin{exercise}
In perfect hashing, suppose the top level sends $n_j = 4$ keys into some slot $j$. What size $m_j$ does the secondary table get? Using the collision-pair argument of Lemma~\ref{ch11:perfect-lemma}, give an upper bound on the expected number of colliding pairs in that secondary table, and explain why this means a collision-free secondary function is easy to find.
\end{exercise}

\begin{exercise}
(Challenge.) Explain the ``hash-flooding'' denial-of-service attack: how can an adversary who knows your fixed hash function force a chained hash table's performance to degrade from $O(1)$ to $O(n)$ per operation? Describe precisely how switching to a randomly chosen universal hash function defeats the attack, referencing Theorem~\ref{ch11:universal-theorem}.
\end{exercise}

\begin{exercise}
(Challenge.) Compare chaining and open addressing on \emph{deletion}. Explain why deletion takes direct $O(1)$ pointer operations in a doubly-linked chained table but requires tombstones in open addressing, and describe one concrete failure that occurs if you delete by simply setting the slot to \textsc{nil}. Then argue why a table full of tombstones can make searches slow even when $n$ is small.
\end{exercise}

\begin{exercise}[Bloom filter dimensioning and zero false negatives]
A storage engine stores $n = 5 \times 10^7$ keys and allocates an array of $m = 5 \times 10^8$ bits ($m/n = 10$).
(a) Compute the optimal number of hash functions $k^\star$ to minimize false positives, rounding to the nearest integer.
(b) Using that $k^\star$, calculate the theoretical false-positive probability $\eps$.
(c) Explain physically why a query that observes a single $0$ bit at any of its $k$ probed locations provides an absolute guarantee that the key is absent ($0$ false negatives), regardless of how many other keys have been inserted into the array.
\end{exercise}

\begin{exercise}[Count-Min Sketch under collisions]
A network router uses a Count-Min Sketch with depth $d = 3$ and width $w = 5$. Counters $C[i, j]$ are initially zero.
Three packet streams $A$, $B$, and $C$ arrive with true volumes $f(A) = 10$, $f(B) = 4$, and $f(C) = 2$.
The hash functions map the keys as follows:
$h_1(A)=2$, $h_2(A)=4$, $h_3(A)=1$;
$h_1(B)=2$, $h_2(B)=1$, $h_3(B)=5$;
$h_1(C)=3$, $h_2(C)=4$, $h_3(C)=1$.
(a) Trace the counter matrix $C[1..3, 1..5]$ after processing all arrivals.
(b) Compute the estimated frequencies $\hat{f}(A)$, $\hat{f}(B)$, and $\hat{f}(C)$ using the minimum rule $\hat{f}(x) = \min_{1 \le i \le 3} C[i, h_i(x)]$.
(c) Verify that $\hat{f}(x) \ge f(x)$ for all three keys, and identify which row provided the exact count for each key. Why can hash collisions never cause undercounting?
\end{exercise}
